{\large\textbf{Immigration Data:}} 
The Immigration data is sourced from \cite{ImmigrationInnovationGrowth}. They present three variables of interest: Immigration, Non-European Immigration, and an Immigration instrument for every five-year period from 1975 to 2010 at the county-level. Table \ref{tab:immigration} provides summary statistics on these key variables. The immigration instruments were constructed using data from the Integrated Public Use Microdata Series (IPUMS) sample waves from 1880 to 2000 of the US census\footnote{The 1940, 1950, and 1960 waves weren't included as they did not include information on immigration. The 1890 census was lost in a fire. For more information, see \href{https://www.familysearch.org/en/blog/what-happened-to-the-1890-us-census}{\textcolor{accentblue}{here.}}}, along with the 2006 to 2010 five-year sample of the American Community Survey. The instruments predict plausibly exogenous non-European migration to US counties. To further illustrate the spatial and time variation of the data, Figure \ref{fig:immigration_map} shows the immigration shocks for every 5-year period in my sample.

{\large\textbf{Business Level Micro-Data:}} I utilize the Survey of Business Owners from the \cite{sbo}. This is a public-use microdata \footnote{The US Census Bureau was experimenting with providing more public-use data in 2007. Shortly after the release of this dataset, they reviewed their privacy policy and chose to keep future iterations of this survey private. }  cross-sectional sample. The data for this sample were collected along with the 2007 Economic Census. It reports data from about 2.3 million nonfarm businesses and includes population weight. The weighted sample is about 26 million businesses and provides a snapshot of all firms in the United States. The survey collected data on the personal characteristics of up to 4 business owners for each business. The data is reported at the individual business level. For this study, the most important question asked is ``Whether the owner was born in the United States?''. I used this variable to identify Immigrant versus American-owned businesses. I assume that if an owner is born outside the US, they are an immigrant. I construct 3 categories of ownership type:
    \begin{enumerate}[1)]
        \item \textbf{Fully American:} If at least one owner reports being born in the US, and no owner reports being born outside the US.
        \item \textbf{Fully Immigrant:} If at least one owner reports being born outside the US, and no owner reports being born in the US.
        \item \textbf{Mixed}: If at least one owner reports being born in the US and one reports being born outside the US.
    \end{enumerate}
    Table \ref{tab:sbo} reports various firm-level attributes based on the ownership type. For my analysis, I generally limit my sample to just the first two groups. Since the owner's place of birth is not reported for all observations, my restricted sample includes only about 13.4 million weighted observations. I believe this sample still provides ample coverage of the US Economy and allows us to test my model predictions.

{\large\textbf{Business Dynamism:}}  The Business Dynamics Statistics are sourced from the \cite{bds}. Data is reported at the county-year level. The main variables I utilize from this dataset are average employment (firm size) and exit rate. It also includes data stratified by firm age, which I utilize for my heterogeneity analysis. Overall, the coverage spans from 1978 to 2010 and has over 100,000 un-stratified observations.

{\large \textbf{Wage Data:}} I use Quarterly Census of Employment and Wages from the \cite{bls} for my wage data. I utilize the annual average pay variable at the county-level. My data spans from 1990 to 2010 and covers 3,125 US counties across all 50 states.

{\large \textbf{County Population:}} Finally, as a control, I use NBER's Intercensal US County Population Data \citep{countpop}. The dataset covers the universe of US counties from 1970 to 2014.
